\section{Antecedentes, problema y solución propuesta}

\subsection{Descripción de la problemática u oportunidad}

Los cerros orientales de Bogotá conforman un borde urbano-forestal donde la vegetación de montaña colinda de manera directa con barrios, vías y equipamientos de la ciudad, en temporadas secas, se corre el riesgo de presentar incendios en estas zonas, que posteriormente se pueden expandir a las áreas urbanas. La responsabilidad de su control recae en los cuerpos de bomberos, según lo establece la Ley 1575 de 2012, mientras que la coordinación general del riesgo se organiza a través del Sistema Nacional de Gestión del Riesgo de Desastres definido por la Ley 1523 de 2012, el cual opera mediante niveles territoriales que deben agotarse de manera sucesiva antes de escalar a instancias superiores. Esta estructura implica que las decisiones más determinantes durante las primeras horas de un evento se toman en el nivel local, con los recursos y la información disponibles en ese momento.

La entrevista realizada a la coronel Ana Milena Mejía de la Fuerza Aérea Colombiana, permitió establecer que el reporte comunitario es el primer insumo con el que cuentan los organismos de socorro y que este llega en términos descriptivos y aproximados: una ubicación referida a una vereda o a un sector reconocible, además de mecanismos como los observadores comunitarios del riesgo y la cartografía social, que ofrece identificación temprana de los siniestros, pero depende de la disponibilidad de voluntarios y carece de detalles importantes como el comportamiento futuro del fuego.
  
Un aspecto central identificado en la entrevista es que la diferencia entre una atención exitosa y una que se sale de control no radica en el procedimiento sino en la oportunidad. Los protocolos de respuesta se encuentran estandarizados, de modo que lo determinante no es \emph{cómo} actuar sino \emph{qué tan rápido} se puede actuar. Esta afirmación se ilustra con lo ocurrido en Tierra Alta durante la emergencia de Córdoba, que pese a ser un desastre natural diferente, ejemplifica el punto que se quiere abordar. Donde una estimación anticipada y elemental del tiempo de llegada del caudal hubiera permitido al alcalde ordenar una evacuación preventina y salvar un número considerable de vidas, cosa que lastimosamente no se logró. El caso evidencia que una proyección temprana, aun cuando sea aproximada, puede modificar de manera sustancial el curso de las decisiones durante una emergencia.

A nivel tecnológico, las herramientas actualmente empleadas se orientan a representar dónde se encuentra la emergencia y no hacia dónde puede dirigirse. El trabajo se apoya en sistemas de información geográfica como ArcGIS para graficar los puntos de incendio y en registros históricos para construir mapas de calor que orientan la distribución anticipada de recursos, siempre bajo la restricción de un presupuesto limitado. Adicionalmente, la información se encuentra sectorizada: cada entidad opera sus propios sistemas y no existe un consolidado común, situación que la Plataforma Nacional para la Gestión del Riesgo de Desastres busca corregir desde su lanzamiento reciente, aunque su adopción todavía se encuentra en proceso. A esto se suman experiencias previas de desarrollo tecnológico que no alcanzaron los resultados esperados, como la plataforma Estrategos de la Defensa Civil Colombiana, cuya inversión no logró sustituir ni integrar los sistemas existentes. Este antecedente muestra que las soluciones concebidas como plataformas institucionales de gran alcance tienden a enfrentar dificultades de adopción, sostenimiento y articulación.

Finalmente, la conectividad constituye una restricción que no puede pasarse por alto. Durante el terremoto del 10 de agosto de 2026 los sistemas de comunicación colapsaron por completo y no fue posible establecer contacto con Quibdó, lo que obligó a desplazar antenas satelitales y repetidores para restablecer un canal mínimo de información. Los respaldos disponibles, como las redes de radio institucionales, no tienen cobertura uniforme en todo el territorio. En consecuencia, cualquier herramienta que dependa de procesamiento intensivo en línea durante el evento resulta frágil justamente en el momento en que más se necesita.

De lo anterior se desprende una oportunidad clara para la Ingeniería de Sistemas. Existe información suficiente sobre el territorio, condiciones climáticas y comportamiento del fuego para anticipar escenarios, pero esa información no se encuentra dispuesta en un formato que los organismos de respuesta puedan consultar de manera inmediata y comprensible durante los primeros minutos de un incendio. Preparar con antelación esa estimación, en lugar de calcularla bajo presión y con conectividad incierta, permitiría reducir el tiempo entre la detección del evento y la decisión operativa, que es precisamente donde se concentra la mayor posibilidad de contener un incendio forestal antes de que alcance una magnitud mayor.


  
\subsection{Formulación del problema}

Durante la atención de incendios forestales en zonas de interfaz urbano-forestal (WUI) y áreas de alta complejidad geográfica como los Cerros Orientales de Bogotá, los organismos de respuesta enfrentan deficiencias críticas derivadas de la reacción tardía y de la ausencia de herramientas que anticipen tempranamente la propagación del fuego. Tal como se ha identificado en la práctica institucional, la capacidad de mitigar un evento depende estrechamente de la oportunidad en la información y de la velocidad de reacción en los primeros minutos de la emergencia. 

No obstante, los modelos convencionales de simulación del comportamiento del fuego exigen un alto costo computacional en tiempo real que resulta inviable durante una emergencia activa. A esto se suman las recurrentes fallas en las redes de comunicación y conectividad móvil como las experimentadas en situaciones de crisis nacional, las cuales imposibilitan la ejecución de procesos de cálculo complejos o la consulta de datos en línea sobre el terreno. 

En consecuencia, surge el siguiente interrogante de investigación: ¿Cómo proporcionar de manera oportuna a los cuerpos de bomberos y organismos de socorro una estimación anticipada de la propagación de incendios forestales en las etapas iniciales en los cerros orientales de Bogotá, superando las restricciones de conectividad y el elevado costo computacional en tiempo real, para optimizar el apoyo a la toma de decisiones operativas?

  

\subsection{Propuesta de solución}

Se propone desarrollar una aplicación web o móvil de apoyo a la toma de decisiones frente a incendios forestales en los Cerros Orientales de Bogotá. Su propósito será ofrecer a los bomberos una visualización del comportamiento probable del fuego desde las primeras etapas del evento, mediante una animación o un video elaborado a partir de simulaciones realizadas previamente. La aplicación no ejecutará una simulación completa durante la emergencia; en cambio, pondrá a disposición información preparada con antelación para que pueda consultarse con mayor oportunidad.

La solución se sustentará en una biblioteca híbrida de escenarios. FARSITE aportará simulaciones de propagación en áreas amplias, mientras que WFDS contribuirá con resultados de mayor detalle en sectores críticos de WUI. Estos escenarios se construirán previamente mediante infraestructura HPC y considerarán información climática y de viento obtenida a través de una API. A partir de los resultados almacenados, la aplicación presentará una representación visual que facilite interpretar la posible evolución del incendio y su relación con el territorio.

\subsection{Justificación de la solución}

La entrevista con la coronel Ana Milena Mejía permitió identificar que la oportunidad de la información es determinante en la atención de incendios forestales. Según su experiencia, la detección temprana y el reporte oportuno aumentan las posibilidades de controlar un evento antes de que alcance una mayor magnitud, especialmente en zonas rurales donde la reacción puede tardar más. También destacó el valor de los observadores comunitarios del riesgo y de la cartografía social para reconocer áreas vulnerables y alertar a las entidades responsables.

Además, las comunicaciones convencionales pueden fallar durante una emergencia, como ocurrió durante el terremoto del 10 de agosto de 2026, cuando se perdió la comunicación con Quibdó. Esta limitación refuerza la decisión de ejecutar y almacenar las simulaciones con antelación, en lugar de depender de una simulación completa en el momento del evento, cuando la conectividad puede ser inestable o inexistente. De este modo, la visualización puede basarse en resultados previamente generados y no requiere iniciar un proceso de cálculo complejo y distribuído durante una situación crítica.
